\section{Data, models and reproducibility details}
\label{app:setup}
All experiments use public binary classification tables. German Credit uses the UCI mixed-type 20-variable version, with original label 2 (bad credit) as positive. Pima uses OpenML id 37, version 1; zero glucose, blood pressure, skinfold thickness, insulin and BMI are treated as missing, with training-median imputation. We neither deduplicate nor remove rows. ACS uses folktables 0.0.12, \texttt{ACSIncome}, California, 2018, 1-Year person records. The package's standard filters require age greater than 16, personal income greater than 100, positive weekly working hours, and person weight at least one; the target is personal income above \$50,000. The ten features are age, class of worker, education, marital status, occupation, birthplace, relationship, usual weekly hours, sex and race. We use the source's numeric category codes for trees, with no population weighting. This is an unweighted record-level computational benchmark, not an ACS population estimate.

The master seed is 20261102. A stratified 75/25 split defines training and test pools. For ACS only, stratified subsampling limits training to 20,000 observations. Median imputation is fitted on training records; logistic numerical variables are standardized and categorical variables one-hot encoded with unknown categories ignored. The entire encoded group is replaced together when masking an original categorical variable. For trees, each original feature occupies one column and category codes have a numerical ordering; this modeling choice is held fixed throughout. It need not be the best categorical-data modeling choice.

The logistic model uses $C=1$ and at most 2,000 solver iterations. Random forest uses 64 trees, maximum depth 6 and minimum leaf size 5. Histogram gradient boosting uses 80 iterations, at most 15 leaves per tree and $\ell_2$ regularization 1. Remaining scikit-learn defaults are pinned by the environment. We do no tuning or feature selection during training. Each cell has one training seed, as the experiment conditions on the resulting model. A dataset-specific RNG selects 32 test cases and 16 training-background rows without replacement; the selections are matched across models. Full original row indices, fitted preprocessing and models are saved in \texttt{data/*.joblib}.

Explanation seeds are generated by \texttt{numpy.random.SeedSequence} from the master seed and dataset, model, case, budget, replicate and method identifiers. Kernel and permutation seeds use distinct method identifiers. Kernel's unregularized and top-10 ablations reuse the same coalition seeds to reduce unnecessary sampling differences. Separate seeds control the background draws, online timing trials and bootstrap. The implementation stores the actual active columns, coalition counts, all warnings, every estimate, and all exact references.

\section{Why the margin implication holds}
\label{app:proof}
For any coordinate, $\big||\widehat\phi_j|-|\phi_j|\big|\leq E$. If $j\in T_3(\phi)$ and $k\notin T_3(\phi)$ with $\Delta_3>0$, then
\begin{align*}
 |\widehat\phi_j|-|\widehat\phi_k|
 &\geq |\phi_j|-|\phi_k|-2E\\
 &\geq \Delta_3-2E.
\end{align*}
This is positive when $E<\Delta_3/2$, so no outside coordinate can enter the top-three set. If $\Delta_3=0$, the strict sufficient condition is empty and cannot certify the arbitrary tie-broken set. Similarly, if $|\phi_j|>E$, an estimate cannot cross zero. These implications require an error bound relative to the target; two agreeing runs do not supply such a bound.

\section{Independent validation of the exact audit}
\label{app:validation}
\texttt{code/validate.py} implements three independent checks before linking experimental outputs to the paper. First, a six-player numerical test compares the subset-weighted formula against summing marginal contributions over all 720 permutations. Second, an additive eight-player game checks that sampled Kernel SHAP recovers its known coefficients. Third, exhaustive eight-player Kernel SHAP is compared with the subset formula for a nonadditive numerical game. These synthetic objects are implementation checks only and do not appear as empirical evidence about public datasets.

A fourth check uses all nine actual dataset/model cells: for each of LR, RF and HGB on each dataset, a real-data KernelExplainer operating on the 16 raw background rows is compared with the binary-mask cached-game wrapper on the same test input, seed, identity link, an $8d$ coalition cap and no regularization. Agreement to numerical precision verifies that grouping, masking, model scale, endpoints and payoff lookup reproduce live evaluation. The machine-readable maximum errors are in \texttt{results/validation.json}; assertions require every difference below $10^{-10}$. Main reference efficiency has maximum residual \ReferenceEfficiency{}. Exhaustive smaller-game Kernel estimates have maximum absolute error \ExhaustiveError{} relative to the reference. We also verify that no observed run violates the deterministic strict-margin implication.

\section{Complete metrics and the meaning of the intervals}
\label{app:metrics}
Table~\ref{tab:full} reports the component metrics underlying the main joint-success table. ``Signs'' is the fraction of runs preserving the signs of all reference contributions with magnitude at least $.01$. ``Repeat top-3'' averages 32 disjoint seed pairs per case. ``Cert. cases'' counts cases with a pointwise one-sided $95\%$ lower bound on joint seed-success probability at least $.95$. It is not a simultaneous certification of all cases. The full per-case CSV additionally reports maximum error, relative $\ell_1$ error, alternative numerical thresholds and actual coalition costs.

Intervals in the main tables resample cases, not random seeds or model/case combinations. At an exhaustive coalition cap the algorithm is deterministic for a fixed target, so a degenerate interval is expected and must not be interpreted as population-wide certainty. With only 32 cases, heterogeneity and rare failure probabilities are imperfectly resolved. We do not infer model-class rankings from predictive scores or from overlapping success-rate intervals.

\begin{table*}[t]
\centering\scriptsize
\begin{tabular}{llrrrrrr}
\toprule
Data/model & Budget & Num. & Top-3 & Joint & Signs & Repeat top-3 & Cert. cases \\
\midrule
German/LR & 2d & 11.1 & 60.2 & 9.8 & 55.4 & 42.2 & 0/32 \\
German/LR & 8d & 97.4 & 93.2 & 90.7 & 100.0 & 91.9 & 18/32 \\
German/LR & 32d & 100.0 & 96.6 & 96.6 & 100.0 & 94.8 & 27/32 \\
German/LR & 128d & 100.0 & 97.9 & 97.9 & 100.0 & 97.1 & 29/32 \\
German/RF & 2d & 24.8 & 52.6 & 20.3 & 78.3 & 33.8 & 0/32 \\
German/RF & 8d & 99.8 & 91.5 & 91.3 & 100.0 & 88.7 & 18/32 \\
German/RF & 32d & 100.0 & 96.2 & 96.2 & 100.0 & 94.1 & 26/32 \\
German/RF & 128d & 100.0 & 98.4 & 98.4 & 100.0 & 97.5 & 28/32 \\
German/HGB & 2d & 0.3 & 31.9 & 0.3 & 28.4 & 14.6 & 0/32 \\
German/HGB & 8d & 55.4 & 89.5 & 48.6 & 96.0 & 83.5 & 2/32 \\
German/HGB & 32d & 99.1 & 96.6 & 95.8 & 100.0 & 95.8 & 22/32 \\
German/HGB & 128d & 100.0 & 98.6 & 98.6 & 100.0 & 98.4 & 31/32 \\
Pima/LR & 2d & 32.2 & 69.5 & 29.8 & 82.7 & 54.7 & 0/32 \\
Pima/LR & 8d & 96.9 & 91.5 & 88.8 & 100.0 & 89.3 & 21/32 \\
Pima/LR & 32d & 100.0 & 100.0 & 100.0 & 100.0 & 100.0 & 32/32 \\
Pima/LR & 128d & 100.0 & 100.0 & 100.0 & 100.0 & 100.0 & 32/32 \\
Pima/RF & 2d & 26.2 & 68.9 & 23.5 & 80.3 & 53.0 & 0/32 \\
Pima/RF & 8d & 96.7 & 95.1 & 92.7 & 100.0 & 92.6 & 18/32 \\
Pima/RF & 32d & 100.0 & 100.0 & 100.0 & 100.0 & 100.0 & 32/32 \\
Pima/RF & 128d & 100.0 & 100.0 & 100.0 & 100.0 & 100.0 & 32/32 \\
Pima/HGB & 2d & 1.3 & 43.2 & 1.1 & 57.3 & 24.0 & 0/32 \\
Pima/HGB & 8d & 45.2 & 86.2 & 40.2 & 95.6 & 78.8 & 0/32 \\
Pima/HGB & 32d & 100.0 & 100.0 & 100.0 & 100.0 & 100.0 & 32/32 \\
Pima/HGB & 128d & 100.0 & 100.0 & 100.0 & 100.0 & 100.0 & 32/32 \\
ACS/LR & 2d & 12.2 & 55.5 & 11.2 & 61.8 & 36.7 & 0/32 \\
ACS/LR & 8d & 78.7 & 93.7 & 74.2 & 98.3 & 90.2 & 9/32 \\
ACS/LR & 32d & 100.0 & 98.3 & 98.3 & 100.0 & 97.5 & 29/32 \\
ACS/LR & 128d & 100.0 & 100.0 & 100.0 & 100.0 & 100.0 & 32/32 \\
ACS/RF & 2d & 18.2 & 65.7 & 16.5 & 77.5 & 48.5 & 0/32 \\
ACS/RF & 8d & 95.9 & 94.9 & 91.1 & 99.9 & 91.9 & 15/32 \\
ACS/RF & 32d & 100.0 & 99.5 & 99.5 & 100.0 & 99.0 & 29/32 \\
ACS/RF & 128d & 100.0 & 100.0 & 100.0 & 100.0 & 100.0 & 32/32 \\
ACS/HGB & 2d & 2.8 & 38.5 & 2.2 & 57.3 & 23.6 & 0/32 \\
ACS/HGB & 8d & 61.2 & 80.2 & 52.1 & 97.8 & 71.7 & 4/32 \\
ACS/HGB & 32d & 99.5 & 91.8 & 91.4 & 100.0 & 88.8 & 16/32 \\
ACS/HGB & 128d & 100.0 & 100.0 & 100.0 & 100.0 & 100.0 & 32/32 \\
\bottomrule
\end{tabular}
\caption{Complete unregularized Kernel SHAP component success rates (percent), with 32 cases per cell and 64 runs per case. The main text reports case-bootstrap intervals for joint success; machine-readable summaries contain intervals for every component.}
\label{tab:full}
\end{table*}

\section{Averaging low-budget explanations}
\label{app:averaging}
Four consecutive independent $2d$ runs are averaged featurewise, yielding 16 independent four-run means per case. We compare their joint success with that of a single $8d$ run, using all 64 $8d$ runs to estimate its success probability. This uses approximately the same interior-payoff cap in the unexhausted regime, but four runs require four endpoint initializations and can repeat coalitions across runs. It is not an exact wall-clock equality: each four-run policy uses four times the realized low-budget counts plus its endpoints; the single-run policy has its own realized count. On small games the cap saturates, which also changes the ratio. We save these costs and interpret the comparison as a budget-allocation diagnostic, not an estimator theorem.

\begin{table*}[t]
\centering\small
\begin{tabular}{llrrr}
\toprule
Data & Method & Four $2d$ runs & One $8d$ run & Difference (pp) \\
\midrule
German & Kernel & 13.9 [11.1, 16.8] & 76.9 [72.0, 81.6] & -62.9 [-66.3, -59.4] \\
German & Permutation & 75.8 [70.5, 80.7] & 73.8 [69.1, 78.2] & 1.9 [-0.1, 4.1] \\
Pima & Kernel & 28.3 [23.0, 33.5] & 73.9 [67.9, 79.4] & -45.6 [-49.3, -41.7] \\
Pima & Permutation & 73.6 [68.6, 78.6] & 75.7 [71.4, 80.0] & -2.1 [-3.9, -0.5] \\
ACS & Kernel & 17.1 [12.4, 22.0] & 72.4 [65.1, 79.4] & -55.3 [-59.5, -51.1] \\
ACS & Permutation & 77.1 [71.2, 82.8] & 77.4 [71.5, 83.1] & -0.3 [-2.5, 1.7] \\
\bottomrule
\end{tabular}
\caption{Joint success (percent, 95\% case-bootstrap intervals), averaged equally over fitted models. Difference is four-run averaging minus a single larger run, with a paired case bootstrap. Endpoints and saturation prevent exact equality of total costs.}
\label{tab:averaging}
\end{table*}

\section{Post hoc numerical-solver control}
\label{app:solver}
Inspection of the poor four-run averaging results revealed extremely large low-budget outputs. We therefore added this diagnostic after the prospective grid was complete; it does not replace the primary results. We replay every $2d$ Kernel run with the same seed, masks, weights and payoffs and verify exact equality with its saved package output. Eliminating the final active feature using the efficiency constraint, we solve $\min_\beta\|W^{1/2}(X\beta-y)\|_2$ directly with NumPy's SVD least squares and default rank tolerance, instead of solving the normal equations. Under rank deficiency this is one coordinate-system-dependent minimum-norm solution, not a newly justified Shapley estimator.

Of \SolverRuns{} replayed designs, \DeficientDesigns{} are rank deficient under the SVD tolerance. The package's $E>1$ extreme errors occur \ExtremePackage{} times and disappear in the SVD control. Table~\ref{tab:solver} shows that low-budget numerical and joint success remain poor: numerical conditioning explains the extreme tail, not the entire finite-budget error. No nonfinite or extreme primary run was discarded or clipped. This post hoc result is implementation- and version-specific, and does not establish an optimal stabilization policy.

\begin{table*}[t]
\centering\small
\begin{tabular}{llrrrr}
\toprule
Data & Solver & Numeric & Top-3 & Joint & $E>1$ \\
\midrule
German & Package & 12.1 [9.5, 14.6] & 48.2 [44.9, 51.6] & 10.1 [7.8, 12.5] & 3.5 [3.1, 4.0] \\
German & SVD & 13.0 [10.2, 15.7] & 53.3 [49.5, 57.2] & 10.9 [8.3, 13.3] & 0.0 [0.0, 0.0] \\
Pima & Package & 19.9 [15.9, 24.1] & 60.5 [56.4, 64.4] & 18.1 [14.2, 22.4] & 0.2 [0.1, 0.3] \\
Pima & SVD & 20.0 [16.0, 24.3] & 60.9 [56.7, 64.9] & 18.2 [14.3, 22.5] & 0.0 [0.0, 0.0] \\
ACS & Package & 11.1 [8.2, 14.2] & 53.2 [49.3, 57.3] & 9.9 [7.2, 12.9] & 0.8 [0.6, 1.0] \\
ACS & SVD & 11.2 [8.3, 14.4] & 54.2 [50.0, 58.4] & 10.1 [7.4, 13.0] & 0.0 [0.0, 0.0] \\
\bottomrule
\end{tabular}
\caption{Post hoc $2d$ solver control with identical sampled payoffs. Percent, with 95\% matched-case bootstrap intervals, averaging the three fixed models in each dataset. $E>1$ denotes maximum attribution error exceeding one probability unit. SVD removes extreme outputs but does not make the lowest budget accurate.}
\label{tab:solver}
\end{table*}

\section{Uncached deployment timing}
\label{app:timing}
Table~\ref{tab:timing} times the real scikit-learn models through an ordinary real-input KernelExplainer, including construction, background evaluation, coalition generation and regression. Each cell is the median of five independently seeded trials on the first selected case, and includes preprocessing for logistic regression. The CPU is Intel Xeon Gold 6230 at nominal 2.10 GHz; the host exposes 40 CPUs, but the container quota is four CPU equivalents. All numerical libraries are limited to one thread per process. These timings were collected alongside the offline audit and are subject to contention. They are descriptive latency measurements, not portable performance guarantees or per-case latency distributions. Actual model-row counts and individual times accompany the code.

Offline exact tables use batches of 2,048 coalitions. For logistic regression, we algebraically combine featurewise logit contributions and apply the sigmoid to compute the same masked probabilities efficiently; trees use their actual probability predictors. Cached payoff tables accelerate repeated explanations without changing their sample access. The retained main enumeration outputs consumed \GamesCpuHours{} process CPU hours; retained enumeration, repeated explanations, background experiments and solver diagnostics together consumed \TotalCpuHours{} process CPU hours. These totals exclude downloads, compilation, failed preparation, timing, analysis and pilots overwritten by a deterministic rerun, and are not claimed as end-to-end accounting. The logs preserve those stages separately. This is well within the task's approximate 24 CPU-hour allowance.

\begin{table*}[t]
\centering\small
\begin{tabular}{llrrrr}
\toprule
Dataset & Model & $2d$ & $8d$ & $32d$ & $128d$ \\
\midrule
German & LR & 105.3 & 114.3 & 290.9 & 698.1 \\
German & RF & 100.9 & 106.4 & 297.2 & 910.2 \\
German & HGB & 90.1 & 100.3 & 307.1 & 1402.5 \\
Pima & LR & 6.7 & 7.8 & 96.7 & 10.6 \\
Pima & RF & 98.2 & 105.0 & 117.7 & 183.1 \\
Pima & HGB & 8.3 & 84.0 & 104.4 & 105.4 \\
ACS & LR & 96.8 & 111.4 & 291.3 & 620.4 \\
ACS & RF & 91.9 & 100.9 & 191.0 & 317.7 \\
ACS & HGB & 8.2 & 13.9 & 113.7 & 303.6 \\
\bottomrule
\end{tabular}
\caption{Median milliseconds for one actual online unregularized KernelExplainer run. Five trials per dataset/model/budget, on one selected input. The budget labels are caps; full smaller games may already be exhausted.}
\label{tab:timing}
\end{table*}

\section{Research conduct and release checklist}
\label{app:checklist}
We use the official ICML 2026 style in its preprint mode. Its author instructions set an eight-page main-text submission limit and require an Impact Statement; references and appendices are outside that limit. Those instructions do not specify a separate NeurIPS-style paper checklist. The following is our supplementary disclosure checklist, not a substituted conference form.
\begin{enumerate}
\item \textbf{Claims and scope: Yes.} The abstract, Sections~\ref{sec:method} and~\ref{sec:discussion} specify the fixed-game, probability-scale, finite-background estimand and exclude causal, conditional-SHAP and model-retraining claims.
\item \textbf{Limitations: Yes.} The discussion covers small test samples, one split and model seed, limited datasets, coded categories, arbitrary display/tolerance choices, implementation dependence and background semantics.
\item \textbf{Theory assumptions and proof: Yes.} The elementary ranking/sign statements are deterministic and proved in Appendix~\ref{app:proof}; no unproved concentration certificate is claimed.
\item \textbf{Code and data access: Yes.} All source code, pinned dependencies, seeds, data URLs/hashes, split indices, frozen models and experimental arrays are included in the workspace; the lightweight source archive omits large raw-data caches. Public data are downloaded without accounts; attribution and license details are in the README.
\item \textbf{Experimental detail and uncertainty: Yes.} The protocol records sampling, costs and metrics before explanation analysis. Case-level bootstrap intervals and pointwise binomial lower bounds are distinguished; seeds are not treated as independent cases.
\item \textbf{Compute and failures: Yes.} Actual CPU output accounting and online latency are reported separately; all singular-regression warnings and the failed pandas preprocessing run are retained. The full planned primary grid and named secondary analyses were completed; the exploratory solver control was added after seeing extreme outputs; larger neural models and alternative modern estimators were outside the design, not failed completed experiments.
\item \textbf{Ethics and human subjects: No new subjects.} We use public historical tables and no human study, contact, annotation labor or private data. The Impact Statement notes representational limitations and risks of overclaiming explanation validity.
\item \textbf{Automated assistance and source independence: Disclosed.} An autonomous agent implemented, ran and wrote this study. Only explicitly permitted starter articles were accessed directly; no topic search, forbidden-paper identity inference or unknown follow-up page was used. Source access and failed URL probes are documented in \texttt{RESEARCH\_LOG.md}. No known accidental encounter with restricted content occurred.
\end{enumerate}
